\documentclass[12pt,a4paper]{article}

\usepackage[utf8]{inputenc}
\usepackage[T1]{fontenc}
\usepackage[english]{babel}
\usepackage{geometry}
\geometry{margin=2.5cm}
\usepackage{amsmath}
\usepackage{booktabs}
\usepackage{float}
\usepackage{parskip}
\usepackage{hyperref}
\usepackage{algorithm}
\usepackage{algpseudocode}

\title{
  \textbf{Assignment \#1 — Graph Traversals}\\[0.4em]
  \large Artificial Intelligence and Knowledge Engineering\\
  SI4024L, 2025/2026
}
\author{Oscar Lopez}
\date{\today}

\begin{document}
\maketitle
\tableofcontents
\newpage

% ─────────────────────────────────────────────────────────────────────────────
\section{Introduction}
% ─────────────────────────────────────────────────────────────────────────────

This assignment implements pathfinding algorithms on a real railway timetable:
the Koleje Dolnośląskie network in Poland, provided in GTFS format.

Two tasks are solved:

\begin{itemize}
  \item \textbf{Task 1:} Find the best route between two stops, minimising
        either travel time (\texttt{t}) or number of transfers (\texttt{p}),
        using Dijkstra and A*.
  \item \textbf{Task 2:} Given a start stop and a list of stops to visit,
        find the shortest closed tour using Tabu Search.
\end{itemize}

The main challenge is that the graph is \textbf{time-dependent}: the next
available train depends on when you arrive at a stop, so the cost of a route
changes throughout the day.

% ─────────────────────────────────────────────────────────────────────────────
\section{Data and Graph}
% ─────────────────────────────────────────────────────────────────────────────

\subsection{GTFS Files}

The GTFS dataset consists of CSV files describing the timetable:

\begin{table}[H]
\centering
\begin{tabular}{ll}
\toprule
\textbf{File} & \textbf{Content} \\
\midrule
\texttt{stops.txt}           & Stop names, GPS coordinates, platform grouping \\
\texttt{routes.txt}          & Train lines (D1, D6, D66, \ldots) \\
\texttt{trips.txt}           & Individual runs of each line \\
\texttt{stop\_times.txt}     & Departure and arrival times per stop per trip \\
\texttt{calendar.txt}        & Which days each service runs \\
\texttt{calendar\_dates.txt} & Exceptions (holidays, extra services) \\
\bottomrule
\end{tabular}
\caption{GTFS files used in this assignment.}
\end{table}

\subsection{Filtering Services by Date}

Before building the graph, only the trips active on the query date are loaded.
A service is active if its weekday flag in \texttt{calendar.txt} matches the
query day, and no exception in \texttt{calendar\_dates.txt} cancels it.
This ensures that Wednesday and Saturday return different routes.

\subsection{Graph Structure}

The graph nodes are stop platforms. Each edge represents one segment of a trip:

\[
  \text{edge} = \langle \text{from\_stop},\; \text{to\_stop},\;
  t_{\text{dep}},\; t_{\text{arr}},\; \text{trip\_id},\; \text{route} \rangle
\]

Edges are stored sorted by departure time, so the next available departure from
any given time can be found in $O(\log n)$ with binary search.

\subsection{Transfer Penalty}

Changing trains at a station requires walking to another platform. A minimum
transfer time of 2 minutes is enforced: a departure on a different train is
only considered if it leaves at least 120 seconds after the passenger's arrival.

% ─────────────────────────────────────────────────────────────────────────────
\section{Task 1: Shortest Path}
% ─────────────────────────────────────────────────────────────────────────────

\subsection{Search State}

The search state is the pair $(\text{stop}, \text{trip})$. Tracking the current
trip is necessary to correctly apply the transfer penalty: staying on the same
train has no penalty, while boarding a different one requires at least 2 minutes.

\subsection{Dijkstra}

Dijkstra explores states in order of increasing cost, always expanding the
cheapest known state next.

\paragraph{Criterion \texttt{t} — travel time.}
The cost of a state is the arrival time in seconds. The algorithm finds the
state at the destination with the earliest arrival, and the travel time is:
\[
  \text{cost} = t_{\text{arrival}} - t_{\text{start}}
\]

\paragraph{Criterion \texttt{p} — transfers.}
The cost is the pair $(\text{transfers},\; t_{\text{arrival}})$, compared
lexicographically: fewer transfers always wins, and arrival time breaks ties.
A transfer is counted each time the trip changes at a stop.

\subsection{A*}

A* improves on Dijkstra by using a heuristic $h(n)$ that estimates the
remaining cost to the goal. Each state is prioritised by:
\[
  f(n) = g(n) + h(n)
\]
where $g(n)$ is the cost accumulated so far. This guides the search towards
the destination and reduces the number of states explored.

\paragraph{Admissibility.}
For A* to guarantee an optimal solution, the heuristic must never
\emph{overestimate} the true remaining cost. This property is called
\textbf{admissibility}.

\paragraph{Heuristic for criterion \texttt{t}.}
\[
  h(n) = \frac{d_{\text{straight line}}(n,\; \text{goal})}{V_{\max}}
\]
The straight-line distance between two GPS coordinates (haversine formula) divided
by $V_{\max} = 50\;\text{m/s}$ ($\approx 180\;\text{km/h}$) gives a lower bound
on the remaining travel time: no train travels faster than $V_{\max}$, so the
actual time is always at least this value. The heuristic is therefore admissible.

\paragraph{Heuristic for criterion \texttt{p}.}
\[
  h(n) = 0
\]
The zero heuristic is always admissible (transfers remaining $\geq 0$). With
this heuristic A* behaves identically to Dijkstra for the transfer criterion,
since no additional guidance is available without pre-computing a transfer
distance matrix.

\subsection{Bidirectional A*}

Bidirectional A* runs two searches simultaneously: one forward from the origin
and one backward from the goal. They stop when they meet in the middle.

The main challenge is that the backward search cannot use real train schedules
(trains do not run backwards). Instead, it uses the sum of edge durations to
estimate how far back it has reached. When the two frontiers meet, a final
forward A* is run from the meeting point to the goal to produce the actual
feasible path.

This approach explores roughly half the nodes of a unidirectional search on
long-distance queries.

\subsection{Algorithm Comparison}

The following command was used to generate the comparison data:

\begin{verbatim}
GTFS_DATE=2026-03-04 python3 solution.py \
    "Wrocław Główny" "Jelenia Góra" t 06:00:00 --compare
\end{verbatim}

\begin{table}[H]
\centering
\begin{tabular}{lrrr}
\toprule
\textbf{Algorithm} & \textbf{Cost} & \textbf{Nodes visited} & \textbf{Time (ms)} \\
\midrule
Dijkstra          & 146.0 min & 684 & 17.5 \\
A*                & 146.0 min & 499 & 16.6 \\
Bidirectional A*  & 146.0 min & 159 & 11.3 \\
\bottomrule
\end{tabular}
\caption{Algorithm comparison on query \textit{Wrocław Główny $\to$ Jelenia Góra},
         criterion \texttt{t}, departure 06:00, Wednesday 2026-03-04.}
\end{table}

All algorithms find the same optimal cost. A* visits fewer nodes than Dijkstra
because the heuristic prunes states that are far from the destination.
Bidirectional A* reduces the search space further.

% ─────────────────────────────────────────────────────────────────────────────
\section{Task 2: TSP with Tabu Search}
% ─────────────────────────────────────────────────────────────────────────────

\subsection{Problem}

Given a starting stop $A$ and stops $L = \{p_1, \ldots, p_n\}$ to visit, find
the ordering of $L$ that minimises the total cost of the round trip
$A \to \cdots \to A$. There are $n!$ possible orderings, so an exhaustive
search is only feasible for very small $n$.

The cost of the tour is computed by chaining A* calls: each segment departs at
the actual arrival time of the previous one. This means the cost is
\textbf{time-dependent} and \textbf{asymmetric}:
visiting Legnica before Jelenia Góra may cost more or less than the reverse
order, because the available trains differ.

\subsection{Initial Solution}

A greedy nearest-neighbour heuristic builds the starting tour: from the current
stop, always go next to the closest unvisited stop (by A* cost). This gives a
reasonable starting point for Tabu Search.

\subsection{Tabu Search (2a — Basic)}

Tabu Search is a local search that avoids revisiting recent solutions by
keeping a \textbf{tabu list} of recently used moves.

\begin{algorithm}[H]
\caption{Tabu Search}
\begin{algorithmic}[1]
\State $s \leftarrow \text{initial solution}$,\quad $s^* \leftarrow s$
\State $\mathcal{T} \leftarrow$ empty tabu list
\For{each iteration}
  \State Generate all 2-opt neighbours of $s$
  \State $s' \leftarrow$ best neighbour whose move is not in $\mathcal{T}$
  \State Add move of $s'$ to $\mathcal{T}$
  \State $s \leftarrow s'$
  \If{cost$(s') <$ cost$(s^*)$} $s^* \leftarrow s'$ \EndIf
\EndFor
\State \Return $s^*$
\end{algorithmic}
\end{algorithm}

The \textbf{2-opt} neighbourhood works by reversing a sub-sequence of the tour.
For $n$ stops there are $\binom{n}{2}$ possible 2-opt moves.

In variant~2a the tabu list grows unboundedly, so no move is ever repeated.

\subsection{Variable Tabu List Size (2b)}

A fixed-size FIFO list of size $T$ is used instead. When the list is full, the
oldest move is removed. The following sizes are compared:

\begin{table}[H]
\centering
\begin{tabular}{ll}
\toprule
\textbf{$T$} & \textbf{Effect} \\
\midrule
$\infty$ (unbounded) & No move is ever repeated \\
$\sqrt{n}$           & Small list, more freedom to explore \\
$n$                  & Balanced \\
$2n$                 & Large list, less risk of cycling \\
\bottomrule
\end{tabular}
\caption{Tabu list sizes compared in variant 2b.}
\end{table}

\subsection{Aspiration Criterion (2c)}

The aspiration criterion allows a tabu move if it produces a new global best:
\[
  \text{allow tabu move} \iff \text{cost}(\text{candidate}) < \text{cost}(s^*)
\]
The reasoning is that if a tabu move leads to a better solution than any seen
before, it cannot be creating a cycle. Blocking it would be overly restrictive.

\subsection{Neighbourhood Sampling (2d)}

Instead of evaluating all $\binom{n}{2}$ neighbours, only a random fraction
$\rho$ is evaluated per iteration. With $\rho = 0.5$, each iteration is twice
as fast, allowing more iterations in the same time at the cost of missing some
neighbours in each step.

\subsection{Variant Comparison}

The following command was used to generate the comparison data:

\begin{verbatim}
GTFS_DATE=2026-03-04 python3 solution.py \
    "Wrocław Główny" "Jelenia Góra;Legnica;Wałbrzych Główny" \
    t 06:00:00 --task2 --tabu-compare
\end{verbatim}

\begin{table}[H]
\centering
\begin{tabular}{lrrrr}
\toprule
\textbf{Configuration} & \textbf{Cost (min)} & \textbf{Iters} & \textbf{Asp.} & \textbf{Time (ms)} \\
\midrule
Unbounded $T$       & 345.0 &   4 & 0 &   330 \\
$T = \sqrt{n} = 2$  & 345.0 & 200 & 0 & 12596 \\
$T = n = 4$         & 345.0 &   4 & 0 &   326 \\
$T = 2n = 8$        & 345.0 &   4 & 0 &   321 \\
$+$ Aspiration      & 345.0 &   4 & 0 &   327 \\
$+$ 50\% sample     & 424.0 &   3 & 0 &   144 \\
\bottomrule
\end{tabular}
\caption{Tabu Search variant comparison on
         \textit{Wrocław Główny $\to$ \{Jelenia Góra, Legnica, Wałbrzych Główny\}},
         criterion \texttt{t}, departure 06:00, Wednesday 2026-03-04.}
\end{table}

% ─────────────────────────────────────────────────────────────────────────────
\section{Conclusions}
% ─────────────────────────────────────────────────────────────────────────────

\begin{itemize}
  \item Tracking \texttt{trip\_id} in the search state is essential to handle
        transfers correctly, including the minimum transfer penalty.

  \item The haversine heuristic for A* is admissible and reduces the number of
        nodes explored compared to Dijkstra, while guaranteeing the same
        optimal solution.

  \item Computing tour costs with time-dependent A* (instead of a static
        distance matrix) is necessary for correct asymmetric TSP costs.

  \item Among the Tabu Search variants, the aspiration criterion has the
        greatest practical impact by recovering solutions only reachable
        through tabu moves.
\end{itemize}

\end{document}
